\section{Commuting parent Hamiltonians}\label{sec:commuting_parent_ham}

A parent Hamiltonian has \emph{commuting} local terms when all translated
interaction projectors commute with each other.

\begin{definition}[Commuting parent Hamiltonian]\label{def:is_commuting_parent_ham}
    \lean{MPSTensor.IsCommutingParentHam}
    \leanok
    \uses{def:local_term_parent}
    The parent Hamiltonian with block length~$L$ on $N$ sites has
    \emph{commuting} local terms if
    $h_i h_j = h_j h_i$ for all $i,j \in \{0,\ldots,N{-}1\}$.
    This records the translated commutativity clause
    $[\tau_j(P_L),P_L]=0$ from
    Ref.~\cite[Definition~3.9]{Cirac2016MPDO_arXiv}; the ground-space spanning
    clause in the same definition is a separate parent-Hamiltonian condition.
\end{definition}

\begin{definition}[Nearest-neighbor commuting parent Hamiltonian]\label{def:is_nncph}
    \lean{MPSTensor.IsNNCPH}
    \leanok
    \uses{def:is_commuting_parent_ham}
    The nearest-neighbor commuting condition is the length-$2$ specialization
    of the preceding commutativity clause. For every
    $i,j\in\{0,\ldots,N{-}1\}$,
    \[
        h_i h_j=h_j h_i.
    \]
    Ref.~\cite[Definition~3.9]{Cirac2016MPDO_arXiv} also requires the
    corresponding parent Hamiltonian to have the prescribed ground-space
    spanning property.
\end{definition}

\begin{definition}[Ground state of a nearest-neighbor commuting parent Hamiltonian]%
    \label{def:is_nncph_ground_state}
    \lean{MPSTensor.IsNNCPHGroundState}
    \leanok
    \uses{def:is_nncph, def:is_frustration_free}
    The MPS vector $V^{(N)}(A)$ is a ground state of a nearest-neighbor
    commuting parent Hamiltonian when the length-two local terms commute and
    annihilate the vector:
    \[
        h_i h_j = h_j h_i,\qquad
        h_i V^{(N)}(A)=0.
    \]
    This is the zero-energy part of the ground-state condition appearing in
    Ref.~\cite[Theorem~3.10(iii)]{Cirac2016MPDO_arXiv}; the full source parent
    Hamiltonian condition also includes the spanning assertion from
    Definition~3.9.
    % Scope restriction recorded in
    % docs/paper-gaps/cpsv16_nncph_ground_state_scope.tex.
\end{definition}

\begin{remark}[Elementary consequences of the commuting definition]
    \lean{MPSTensor.IsNNCPH.isCommutingParentHam}
    \lean{MPSTensor.IsNNCPHGroundState.isNNCPH}
    \lean{MPSTensor.IsNNCPHGroundState.isFrustrationFree}
    \lean{MPSTensor.IsNNCPH.isNNCPHGroundState}
    \lean{MPSTensor.IsCommutingParentHam.symm}
    \lean{MPSTensor.IsCommutingParentHam.ham_comm_localTerm}
    \leanok
    \uses{lem:parent_hamiltonian_ff}
    The nearest-neighbor condition is the length-$2$ specialization of the
    general commuting condition. The ground-state condition contains this
    commuting equation and the frustration-free equation for $V^{(N)}(A)$.
    For $N\ge 2$, Lemma~\ref{lem:parent_hamiltonian_ff} with $L=2$ gives
    \[
        h_i\,V^{(N)}(A)=0
    \]
    for every site $i$.
\end{remark}

\begin{theorem}[RFP implies NNCPH]\label{thm:rfp_implies_nncph}
    \lean{MPSTensor.rfp_implies_nncph}
    \leanok
    \uses{def:is_rfp, def:is_nncph, def:normal}
    If $A$ has nonzero virtual dimension and is a normal renormalization fixed
    point, then for every chain length $N \ge 2$ the parent Hamiltonian with
    $L = 2$ has commuting local terms.
    This implication is cited here from the product-of-entangled-pairs
    description in Appendix~B of
    Theorem~3.10 in~\cite{Cirac2016MPDO_arXiv}.
\end{theorem}

\begin{theorem}[RFP gives a ground state of a NNCPH]%
    \label{thm:rfp_implies_nncph_ground_state}
    \lean{MPSTensor.rfp_implies_nncph_ground_state}
    \leanok
    \uses{thm:rfp_implies_nncph, def:is_nncph_ground_state}
    If $A$ has nonzero virtual dimension and is a normal renormalization fixed
    point, then for every chain length $N \ge 2$ the MPS vector $V^{(N)}(A)$ is
    a ground state of a nearest-neighbor commuting parent Hamiltonian:
    \[
        h_i h_j = h_j h_i,\qquad
        h_i V^{(N)}(A)=0.
    \]
    % Scope restriction recorded in
    % docs/paper-gaps/cpsv16_nncph_ground_state_scope.tex.
\end{theorem}

\begin{proof}
    \leanok
    \uses{thm:rfp_implies_nncph, lem:parent_hamiltonian_ff}
    The first equation is Theorem~\ref{thm:rfp_implies_nncph}. The second is
    Lemma~\ref{lem:parent_hamiltonian_ff} with $L=2$.
\end{proof}

\begin{remark}[Product-of-pairs equations for NNCPH]%
    \label{rem:product_pair_projectors}
    \lean{MPSTensor.productPairWindow}
    \lean{MPSTensor.productPairWindow_one}
    \lean{MPSTensor.productPairState}
    \lean{MPSTensor.productPairState_one}
    \lean{MPSTensor.HasProductPairMPV}
    \lean{MPSTensor.HasProductPairLocalProjectors}
    \lean{MPSTensor.HasProductPairLocalProjectors.isNNCPH}
    \lean{MPSTensor.ProductPairBridge}
    \lean{MPSTensor.ProductPairBridge.localTerm_idempotent}
    \lean{MPSTensor.ProductPairBridge.commuting_twoSite_localTerms}
    \lean{MPSTensor.ProductPairBridge.isNNCPH}
    \lean{MPSTensor.AppendixBStructuralData}
    \lean{MPSTensor.AppendixBStructuralData.exists_ofRFP}
    \lean{MPSTensor.AppendixBStructuralData.ofRFP}
    \lean{MPSTensor.AppendixBStructuralData.twoSiteAmplitude}
    \lean{MPSTensor.AppendixBStructuralData.coreTensor}
    \lean{MPSTensor.AppendixBStructuralData.coreTensor_apply}
    \lean{MPSTensor.AppendixBStructuralData.gaugeEquiv_coreTensor}
    \lean{MPSTensor.AppendixBStructuralData.mpv_eq_coreTensor}
    \lean{MPSTensor.AppendixBStructuralData.twoSiteAmplitude_eq_mpv}
    \lean{MPSTensor.AppendixBStructuralData.twoSiteAmplitude_eq_coreTensor_mpv}
    \lean{MPSTensor.AppendixBStructuralData.mpv_eq_productPairState_one}
    \lean{MPSTensor.AppendixBProductPairExtraction}
    \lean{MPSTensor.AppendixBProductPairExtraction.ofCoreTensorFactorization}
    \lean{MPSTensor.AppendixBProductPairExtraction.toProductPairBridge}
    \lean{MPSTensor.AppendixBProductPairExtraction.commuting_twoSite_localTerms}
    \lean{MPSTensor.commuting_twoSite_localTerms_of_rfp_of_appendixBExtraction}
    \lean{MPSTensor.rfp_implies_nncph_of_appendixBExtraction}
    \leanok
    The Appendix~B structural form supplies
    \[
        A^i=X\Lambda U^iX^{-1}.
    \]
    The change of basis does not alter periodic MPS coefficients, so the
    corresponding two-site amplitude is
    \[
        \phi_2(a,b)
        =
        V^{(2)}(A)(a,b)
        =
        V^{(2)}(\Lambda U)(a,b).
    \]
    The missing product-of-pairs input is the even-chain factorization
    \[
        V^{(2n)}(A)(\sigma)
        =
        \prod_{p=0}^{n-1}
        \phi_2\bigl(\sigma_{2p},\sigma_{2p+1}\bigr).
    \]
    The nearest-neighbor parent projectors must also be identified with
    idempotents $p_i$ satisfying
    \[
        h_i(A,2)=p_i,\qquad p_i^2=p_i,\qquad p_ip_j=p_jp_i.
    \]
    These equations give
    \[
        h_i(A,2)h_j(A,2)=h_j(A,2)h_i(A,2),
    \]
    which is the NNCPH conclusion for every finite chain. Thus the remaining
    mathematical input after the Appendix~B structural form is exactly the
    factorization above together with the two-site projector identification.
\end{remark}

\begin{theorem}[Pairwise length-two commutativity implies a renormalization fixed point]\label{thm:nncph_implies_rfp}
    \lean{MPSTensor.nncph_implies_rfp}
    \leanok
    \uses{def:is_rfp, def:is_nncph, def:normal}
    If $A$ is normal, in left-canonical form, and for every chain length
    $N \ge 2$ and every $0\le i,j<N$ its length-two parent interaction
    satisfies the pairwise commutativity equation
    \[
        h_i(A,2)h_j(A,2)=h_j(A,2)h_i(A,2)
    \]
    then $A$ is a renormalization fixed point.
    % Scope restriction: this theorem records only the translated two-site
    % commutativity condition. The source definition's ground-space spanning
    % clause is absent; see docs/paper-gaps/cpsv16_nncph_ground_state_scope.tex.
\end{theorem}
